\documentclass[11pt, oneside]{article}
\usepackage{geometry}
\geometry{letterpaper, margin=1in}
\usepackage[parfill]{parskip}
\usepackage{graphicx}
\usepackage{amssymb, amsmath}
\usepackage{array}
\usepackage{enumitem}
\usepackage[dvipsnames]{xcolor}
\usepackage[most]{tcolorbox}
\usepackage{needspace}

% ---- problem header ----
\newcommand{\problem}[1]{%
  \vspace{1.6em}
  \noindent\rule{\textwidth}{0.8pt}\\[0.3em]
  \noindent{\bf\large Problem #1}\\[0.3em]
  \noindent\rule{\textwidth}{0.4pt}
  \vspace{0.4em}
}

% ---- part label, for multi-part problems ----
\newcommand{\prt}[1]{\Needspace*{6\baselineskip}\vspace{1.1em}\noindent{\bf\large (#1)}\ }

% ---- shaded, rounded solution box ----
\newtcolorbox{sol}[1][Solution]{%
  enhanced, breakable,
  colback=RoyalBlue!4,
  colframe=RoyalBlue!55!black,
  boxrule=0.7pt,
  arc=3.5mm,
  left=4mm, right=4mm, bottom=3mm, top=4.5mm,
  coltitle=white,
  colbacktitle=RoyalBlue!55!black,
  fonttitle=\bfseries\small,
  attach boxed title to top left={xshift=5mm, yshift=-2.6mm},
  boxed title style={arc=1.8mm, boxrule=0pt, left=2.5mm, right=2.5mm},
  title={#1},
  before upper={\setlength{\parskip}{0.55\baselineskip}},
  before skip=10pt, after skip=10pt
}

% ---- inner callout, for warnings inside a solution ----
\newtcolorbox{note}[1][Careful!]{%
  colback=BurntOrange!7,
  colframe=BurntOrange!75!black,
  boxrule=0.6pt, arc=2.5mm,
  left=3mm, right=3mm, top=2mm, bottom=2mm,
  fonttitle=\bfseries\small,
  coltitle=BurntOrange!30!black,
  colbacktitle=BurntOrange!18,
  title={#1},
  before upper={\setlength{\parskip}{0.55\baselineskip}},
  before skip=8pt, after skip=8pt
}

\newcommand{\ans}[1]{\fbox{#1}}

\title{Math F113X: Homework Set 2 --- Solutions}
\author{}
\date{}

\begin{document}
\maketitle
\vspace{-3em}

\begin{center}
\fbox{\parbox{0.9\textwidth}{\centering
Voting Theory (p.~53): \# 12, 13--16, 18, 24ab, and Problem A\\[0.3em]
Weighted Voting (p.~70): \# 2, 4, 5, 6}}
\end{center}

\begin{center}
{\Large\bf Part I: Voting Theory}
\end{center}

%%%%%%%%%%%%%%%%%%%%%%%%%%%%%%%%%%%%%%%%%%%%%%%%%%%%%%%%%%%%
\problem{12}

The downtown business association is electing a new chairperson, and decides to use
approval voting. The tally is below, where each column shows the number of voters
with the particular approval vote. Which candidate wins under approval voting?

\begin{center}
\begin{tabular}{|r|c|c|c|c|c|c|c|}
\hline
\textbf{Number of voters} & 8 & 7 & 6 & 3 & 4 & 2 & 5 \\
\hline
\textbf{A} & X & X & & & X & & X \\
\textbf{B} & X & & X & X & & & X \\
\textbf{C} & & X & X & X & & X & \\
\textbf{D} & X & & X & & X & X & \\
\hline
\end{tabular}
\end{center}

\begin{sol}
Under approval voting, each candidate's score is simply the total number of voters who
marked an X in that candidate's row. Add across each row.

\begin{center}
\begin{tabular}{|l|l|l|}
\hline
\textbf{Candidate} & \textbf{\# of approvals} & \\
\hline
A & $8 + 7 + 4 + 5$ & $= 24$ \\
B & $8 + 6 + 3 + 5$ & $= 22$ \\
C & $7 + 6 + 3 + 2$ & $= 18$ \\
D & $8 + 6 + 4 + 2$ & $= 20$ \\
\hline
\end{tabular}
\end{center}

Candidate A has the most approvals, so the winner is \ans{Candidate A}.
\end{sol}

%%%%%%%%%%%%%%%%%%%%%%%%%%%%%%%%%%%%%%%%%%%%%%%%%%%%%%%%%%%%
\problem{13--16}

\noindent\emph{Problems 13--16 all refer to the same situation:}

\begin{quote}
An election resulted in Candidate A winning, with Candidate B coming in a close
second, and candidate C being a distant third.
\end{quote}

\prt{13} If for some reason the election had to be held again and C decided to
drop out of the election, which caused B to become the winner, which is the primary
fairness criterion violated in this election?

\begin{sol}
This violates the \ans{Independence of Irrelevant Alternatives Criterion (IIA)},
because candidate C is irrelevant to the competition as a ``distant'' third. Removing a
candidate who never had a chance of winning should not change which of the remaining
candidates wins.
\end{sol}

\prt{14} If for some reason the election had to be held again and many people
who had voted for C switched their preferences to favor A, which caused B to become the
winner, which is the primary fairness criterion violated in this election?

\begin{sol}
This violates the \ans{Monotonicity Criterion}. If A gains more support, that should not
decrease A's chance of winning --- yet here A gained support and lost the election.
\end{sol}

\prt{15} If in a head-to-head comparison a majority of people prefer B to A or
C, which is the primary fairness criterion violated in this election?

\begin{sol}
This violates the \ans{Condorcet Criterion}, since B is preferred in a head-to-head
comparison with all other candidates. A Condorcet candidate exists (namely B), but a
different candidate (A) won.
\end{sol}

\prt{16} If B had received a majority of first place votes, which is the primary
fairness criterion violated in this election?

\begin{sol}
This violates the \ans{Majority Criterion} --- a candidate receiving a majority of the
first-place votes should be the winner, but B did not win.
\end{sol}

%%%%%%%%%%%%%%%%%%%%%%%%%%%%%%%%%%%%%%%%%%%%%%%%%%%%%%%%%%%%
\problem{18}

In the election shown below under the Borda Count method, explain why voters in the
second column might be inclined to vote insincerely. How could it affect the outcome
of the election?

\begin{center}
\begin{tabular}{|r|c|c|}
\hline
\textbf{Number of voters} & 20 & 18 \\
\hline
\textbf{1st choice} & A & B \\
\textbf{2nd choice} & B & A \\
\textbf{3rd choice} & C & C \\
\hline
\end{tabular}
\end{center}

\begin{sol}
Recall that under the Borda Count, \emph{all} rankings are weighted and counted, not
just first place. With 3 candidates, a ballot awards 3 points for 1st, 2 for 2nd, and
1 for 3rd. Thus the second-column voters could lower the Borda Count for candidate A by
insincerely reversing their ranks of A and C.

Concretely, here is the count as the ballots stand:

\begin{center}
\begin{tabular}{|l|l|l|}
\hline
\textbf{Candidate} & \textbf{Sincere Borda Count} & \\
\hline
A & $3(20) + 2(18) + 1(0) = 60 + 36$ & $= 96$ \quad $\leftarrow$ \textbf{A wins} \\
B & $3(18) + 2(20) + 1(0) = 54 + 40$ & $= 94$ \\
C & $3(0) + 2(0) + 1(20+18)$        & $= 38$ \\
\hline
\end{tabular}
\end{center}

Now suppose the 18 voters in the second column insincerely rank C above A, so their
ballot reads B, C, A:

\begin{center}
\begin{tabular}{|l|l|l|}
\hline
\textbf{Candidate} & \textbf{Insincere Borda Count} & \\
\hline
A & $3(20) + 2(0) + 1(18) = 60 + 18$ & $= 78$ \\
B & $3(18) + 2(20) + 1(0) = 54 + 40$ & $= 94$ \quad $\leftarrow$ \textbf{now B wins} \\
C & $3(0) + 2(18) + 1(20) = 36 + 20$ & $= 56$ \\
\hline
\end{tabular}
\end{center}

By burying A in last place, the second-column voters drop A's total from 96 to 78 while
leaving B's total unchanged at 94. The winner \ans{switches from Candidate A to
Candidate B}, which is the outcome those voters prefer.
\end{sol}

%%%%%%%%%%%%%%%%%%%%%%%%%%%%%%%%%%%%%%%%%%%%%%%%%%%%%%%%%%%%
\problem{24 (parts a and b)}

Sequential Pairwise voting is a method not commonly used for political elections, but
sometimes used for shopping and games of pool. In this method, the choices are assigned
an order of comparison, called an agenda. The first two choices are compared. The winner
is then compared to the next choice on the agenda, and this continues until all choices
have been compared against the winner of the previous comparison.

\begin{center}
\begin{tabular}{|r|c|c|c|}
\hline
\textbf{Number of voters} & 10 & 15 & 12 \\
\hline
\textbf{1st choice} & C & A & B \\
\textbf{2nd choice} & A & B & D \\
\textbf{3rd choice} & B & D & C \\
\textbf{4th choice} & D & C & A \\
\hline
\end{tabular}
\end{center}

\prt{a} Using the preference schedule above, apply Sequential Pairwise voting to
determine the winner, using the agenda: A, B, C, D.

\begin{sol}
We work down the agenda one comparison at a time.

\vspace{0.5em}
\noindent\textbf{Step 1: Compare A and B.}
\begin{center}
\begin{tabular}{|c|l|}
\hline
\textbf{Candidate} & \textbf{\# of voters preferring} \\
\hline
A & $10 + 15 = 25$ \\
B & $12$ \\
\hline
\end{tabular}
\end{center}
A is preferred, so A advances.

\vspace{0.5em}
\noindent\textbf{Step 2: Compare A and C.}
\begin{center}
\begin{tabular}{|c|l|}
\hline
\textbf{Candidate} & \textbf{\# of voters preferring} \\
\hline
A & $15$ \\
C & $10 + 12 = 22$ \\
\hline
\end{tabular}
\end{center}
C is preferred, so C advances.

\vspace{0.5em}
\noindent\textbf{Step 3: Compare C and D.}
\begin{center}
\begin{tabular}{|c|l|}
\hline
\textbf{Candidate} & \textbf{\# of voters preferring} \\
\hline
C & $10$ \\
D & $15 + 12 = 27$ \\
\hline
\end{tabular}
\end{center}
D is preferred.

\vspace{0.5em}
\noindent The agenda is exhausted, so the winner is \ans{Candidate D}.
\end{sol}

\prt{b} Show that Sequential Pairwise voting can violate the Pareto criterion.

\vspace{0.3em}
\noindent\emph{Pareto Criterion: If \textbf{every} voter prefers candidate A to
candidate B, then B should not be the winner.}

\begin{sol}
Look at how each group of voters ranks B relative to D:

\begin{center}
\begin{tabular}{|r|c|c|c|}
\hline
\textbf{Number of voters} & 10 & 15 & 12 \\
\hline
\textbf{Rank of B} & 3rd & 2nd & 1st \\
\textbf{Rank of D} & 4th & 3rd & 2nd \\
\hline
\end{tabular}
\end{center}

In every single column, B is ranked above D. So \emph{every} voter --- all 37 of them
--- prefers candidate B to candidate D. Yet by part (a), D is the winner.

Since every voter prefers B to D and D still wins, \ans{the Pareto criterion is violated}.
\end{sol}

%%%%%%%%%%%%%%%%%%%%%%%%%%%%%%%%%%%%%%%%%%%%%%%%%%%%%%%%%%%%
\problem{A}

\prt{1} Explain Arrow's Impossibility Theorem in your own words in such a way that a
typical adult could understand it.

\begin{sol}
No matter what voting scheme is used, there may occur outcomes that don't feel fair.

More fully: mathematicians have written down a short list of properties that any
reasonable voting method ``ought'' to have --- for example, if every voter prefers one
candidate to another, the second one shouldn't win, and dropping a hopeless candidate
from the race shouldn't change who wins. Arrow's theorem says that when there are three
or more candidates, no ranked voting method can satisfy all of those properties at once.
Every method has to give something up. So the choice of a voting system is not a matter
of finding the one perfect method --- there isn't one --- but of deciding which kinds of
unfairness we are most willing to live with.
\end{sol}

\prt{2} The President of the United States is elected by the Electoral College. Now
that you have learned a lot about voting schemes, what scheme would \textbf{you} choose
to elect the President of the United States and why?

\begin{sol}[Solution --- many answers are acceptable]
Full credit requires naming a specific scheme and giving reasons that draw on the
properties we studied. A sample response:

\begin{quote}
I would use Instant Runoff Voting with a national popular vote. Under the current
system, a third-party candidate can act as a spoiler --- as in \#20 of Homework 1,
where candidate C's entry flipped the winner from B to A. IRV largely removes that
problem, because if my favorite candidate is eliminated my vote transfers to my next
choice, so I never have to vote insincerely for the ``lesser evil.'' It also guarantees
the winner has majority support after transfers, which plurality does not. I know IRV
is not perfect --- it can violate monotonicity, as we saw --- but I would rather accept
that rare failure than keep a system where voting sincerely can hurt the candidate you
most want to win.
\end{quote}

Other well-argued answers --- keeping the Electoral College, Borda Count, approval
voting, a Condorcet method --- earn full credit if the reasoning is specific and uses
the vocabulary from this unit.
\end{sol}

\prt{3} How has learning about voting schemes changed your view of voting and elections?

\begin{sol}[Solution --- many answers are acceptable]
Full credit requires a specific, reflective response that connects to something actually
studied in this unit. A sample response:

\begin{quote}
I used to think that counting votes was the easy, objective part of an election and
that any disagreement was about the politics. What surprised me is that the counting
method itself decides the outcome: in \#6 of Homework 1, the very same 47 ballots
produced Standard A under plurality, Standard B under Borda Count, Standard A under
IRV, and a tie between B and D under Copeland. Nobody changed their mind --- only the
rule changed. That makes me pay much more attention to \emph{how} an election is run,
and it makes me more sympathetic to arguments about election reform than I was before.
\end{quote}
\end{sol}

%%%%%%%%%%%%%%%%%%%%%%%%%%%%%%%%%%%%%%%%%%%%%%%%%%%%%%%%%%%%
\clearpage
\begin{center}
{\Large\bf Part II: Weighted Voting}
\end{center}

%%%%%%%%%%%%%%%%%%%%%%%%%%%%%%%%%%%%%%%%%%%%%%%%%%%%%%%%%%%%
\problem{2}

Consider the weighted voting system $[31: 10, 10, 8, 7, 6, 4, 1, 1]$.

\prt{a} How many players are there?

\begin{sol}
The numbers after the colon are the players' weights, and there are 8 of them.
So there are \ans{8 players}.
\end{sol}

\prt{b} What is the total number (weight) of votes?

\begin{sol}
Add the weights:
\[
10 + 10 + 8 + 7 + 6 + 4 + 1 + 1 = 47.
\]
The total weight is \ans{47 votes}.
\end{sol}

\prt{c} What is the quota in this system?

\begin{sol}
The quota is the number before the colon: \ans{$q = 31$}. For context, this is
$\tfrac{31}{47} \approx 0.6596$, or about $66\%$ of the total weight.
\end{sol}

%%%%%%%%%%%%%%%%%%%%%%%%%%%%%%%%%%%%%%%%%%%%%%%%%%%%%%%%%%%%
\problem{4}

Consider the weighted voting system $[q: 10, 9, 8, 8, 8, 6]$.

\prt{a} What is the smallest value that the quota $q$ can take?

\begin{sol}
First find the total weight of votes:
\[
10 + 9 + 8 + 8 + 8 + 6 = 49.
\]
The quota must be more than half the total weight; otherwise two opposing
coalitions could both reach the quota at the same time. Since
\[
\tfrac{1}{2}(49) = 24.5,
\]
we need $q \geq 25$, so the smallest quota is \ans{$q = 25$}.
\end{sol}

\prt{b} What is the largest value that the quota $q$ can take?

\begin{sol}
The quota cannot exceed the total weight of 49 --- otherwise no coalition, not even
all the players together, could ever pass a motion. So the largest quota is
\ans{$q = 49$}.
\end{sol}

\prt{c} What is the value of the quota if at least two-thirds of the votes are
required to pass a motion?

\begin{sol}
Two-thirds of the total weight is
\[
\tfrac{2}{3}(49) = 32.\overline{6}.
\]
Since ``at least two-thirds'' must be met with a whole number of votes, we round up:
the quota would be \ans{$q = 33$}.
\end{sol}

%%%%%%%%%%%%%%%%%%%%%%%%%%%%%%%%%%%%%%%%%%%%%%%%%%%%%%%%%%%%
\problem{5}

Consider the weighted voting system $[13: 13, 6, 4, 2]$.

\prt{a} Identify the dictators, if any.

\begin{sol}
Player 1 (call this player $P_1$) has weight equal to the quota, $13 = 13$.
So $P_1$ can pass any motion alone, with no help from anyone: \ans{$P_1$ is a dictator}.
\end{sol}

\prt{b} Identify players with veto power, if any.

\begin{sol}
\ans{$P_1$ has veto power}, since a dictator automatically has veto power ---
the remaining players have combined weight $6 + 4 + 2 = 12 < 13$, so nothing passes
without $P_1$. No one else has veto power.
\end{sol}

\prt{c} Identify dummies, if any.

\begin{sol}
\ans{Players $P_2$, $P_3$, and $P_4$ are dummies.} If there is a dictator, all
other players are dummies: every winning coalition already contains $P_1$, and removing
any of $P_2$, $P_3$, $P_4$ leaves a coalition containing $P_1$, which still wins. So none
of them is ever critical.
\end{sol}

%%%%%%%%%%%%%%%%%%%%%%%%%%%%%%%%%%%%%%%%%%%%%%%%%%%%%%%%%%%%
\problem{6}

Consider the weighted voting system $[11: 9, 6, 3, 1]$.

\prt{a} Identify the dictators, if any.

\begin{sol}
The total weight here is $9 + 6 + 3 + 1 = 19$, and the quota is $11$.

\ans{There is no dictator}, since all the weights are below the quota:
$9 < 11$, $6 < 11$, $3 < 11$, $1 < 11$. No player can pass a motion alone.
\end{sol}

\prt{b} Identify players with veto power, if any.

\begin{sol}
\ans{$P_1$ has veto power.} The players other than $P_1$ have combined weight
\[
6 + 3 + 1 = 10 < 11,
\]
so no motion can pass without $P_1$.

No one else has veto power. For instance, $\{P_1, P_3, P_4\}$ has weight
$9 + 3 + 1 = 13 \geq 11$ and wins without $P_2$, so $P_2$ has no veto; and
$\{P_1, P_2\}$ has weight $15 \geq 11$ and wins without either $P_3$ or $P_4$.
\end{sol}

\prt{c} Identify dummies, if any.

\begin{sol}
\ans{$P_4$ is a dummy.} There is no winning coalition for which $P_4$ is
critical. To see this, list the winning coalitions containing $P_4$ and check whether
removing $P_4$ breaks them:

\begin{center}
\begin{tabular}{|l|c|l|}
\hline
\textbf{Winning coalition} & \textbf{Weight} & \textbf{Weight without $P_4$} \\
\hline
$\{P_1, P_2, P_4\}$      & $16$ & $15 \geq 11$ --- still wins \\
$\{P_1, P_3, P_4\}$      & $13$ & $12 \geq 11$ --- still wins \\
$\{P_1, P_2, P_3, P_4\}$ & $19$ & $18 \geq 11$ --- still wins \\
\hline
\end{tabular}
\end{center}

$P_4$ is never critical, so $P_4$ is a dummy.

By contrast, $P_3$ is \emph{not} a dummy: the coalition $\{P_1, P_3\}$ has weight
$12 \geq 11$ and wins, but without $P_3$ the weight drops to $9 < 11$, so $P_3$ is
critical there.
\end{sol}

\end{document}
